\documentclass[11pt]{article}
\usepackage[utf8]{inputenc}
\usepackage[T1]{fontenc}
\usepackage[spanish,es-tabla]{babel}
\usepackage[margin=2.5cm]{geometry}
\usepackage{amsmath,amssymb}
\usepackage{graphicx}
\usepackage{enumitem}
\usepackage{booktabs}

\pagestyle{plain}
\setlength{\parindent}{0pt}
\setlength{\parskip}{0.5em}

\newenvironment{supuestos}{%
  \noindent\textbf{Supuestos:}
  \begin{itemize}[leftmargin=1.4em, itemsep=1pt]
}{%
  \end{itemize}
}

\newenvironment{dcl}{%
  \noindent\textbf{Diagrama de cuerpo libre --- fuerzas presentes:}
  \begin{enumerate}[leftmargin=1.6em, itemsep=2pt, label=\arabic*.]
}{%
  \end{enumerate}
}

\begin{document}
\shorthandoff{<>}

% ---------- Encabezado ----------
\noindent
\begin{minipage}[t]{0.45\textwidth}
\raggedright
Mec\'anica de Fluidos CI3162\\
Primavera 2026
\end{minipage}%
\hfill
\begin{minipage}[t]{0.45\textwidth}
\raggedleft
Prof.\ Alberto de la Fuente Stranger\\
Auxs.\hspace{0.6cm}Josefina Moreno Faune\\
Sebasti\'an Fredes Ojeda
\end{minipage}

\vspace{3pt}
\hrule
\vspace{-5pt}

\begin{center}
{\Huge\bfseries Pauta oficial explicada --- Control 1}\\[6pt]
10 y 11 de Septiembre 2026
\end{center}

\vspace{10pt}
\noindent\textit{Este documento transcribe y explica paso a paso la pauta manuscrita del
profesor, verificando cada resultado num\'ericamente. Se profundiza especialmente en la
integraci\'on de las Preguntas 3--4 y en el an\'alisis de fuerzas de las Preguntas 7--8,
que son las que generaron dudas. Se usa $g=9{,}81\ \text{m/s}^2$ y $1\ \text{atm}=101\,325$
Pa en todos los c\'alculos.}
\vspace{4pt}
\hrule

%======================================================================
\newpage
\noindent{\LARGE\bfseries Pregunta 1 --- Viscosidad por regresi\'on lineal}

\noindent Una placa de \'area $A$ se desliza a velocidad $U$ sobre una capa de lubricante
de espesor $e$ y viscosidad $\mu$. La fuerza necesaria es, por la ley de Newton de la
viscosidad (esfuerzo de corte uniforme en la capa, $\tau=\mu U/e$):
\begin{equation}
F=\mu\,\frac{U}{e}\,A=\underbrace{\left(\frac{\mu A}{e}\right)}_{m}\cdot U
\end{equation}
Es decir, $F$ es \textbf{lineal} en $U$, con pendiente $m=\mu A/e$: por eso los datos
experimentales (pares $U_i,F_i$) se ajustan con la f\'ormula de regresi\'on lineal dada
en el enunciado, tratando $x_i=U_i$, $y_i=F_i$.

\begin{equation}
m=\frac{\sum(x_i-\bar x)(y_i-\bar y)}{\sum(x_i-\bar x)^2}
\end{equation}
Con los 4 pares de datos de la Tabla 1 ($\bar x=1{,}1325$, $\bar y=1{,}29$):
\begin{center}
\fbox{$m=1{,}0482$}
\end{center}
Despejando $\mu$ de $m=\mu A/e$, con $A=131\ \text{cm}^2=0{,}0131\ \text{m}^2$ y
$e=0{,}06\ \text{mm}=6\times10^{-5}$ m:
\begin{equation}
\boxed{\mu=\frac{m\,e}{A}=\frac{1{,}0482\times6\times10^{-5}}{0{,}0131}\approx0{,}0048\ \text{Pa\,s}}
\end{equation}

%======================================================================
\newpage
\noindent{\LARGE\bfseries Pregunta 2 --- \textquestiondown Se rompe el tubo?}

\noindent Un tubo cerrado en un extremo, lleno de agua hasta $h$, con un pist\'on de masa
$M$ en el extremo abierto, se invierte (el extremo cerrado queda arriba).

\begin{supuestos}
\item El pist\'on est\'a en equilibrio est\'atico: la presi\'on atmosf\'erica desde abajo
sostiene su peso m\'as la presi\'on del agua que empuja hacia abajo.
\item Presi\'on hidrost\'atica est\'andar dentro del tubo (agua en reposo).
\end{supuestos}

\begin{dcl}
\item \textbf{Peso propio $Mg$} --- hacia abajo, siempre presente.
\item \textbf{Presi\'on atmosf\'erica $P_{atm}$ sobre la cara inferior del pist\'on}
--- hacia \textbf{arriba} (empuja el pist\'on hacia adentro del tubo); el pist\'on est\'a
en el extremo abierto, as\'i que esta cara mira directamente a la atm\'osfera.
\item \textbf{Presi\'on del agua $P_b$ sobre la cara superior del pist\'on} --- hacia
\textbf{abajo}; es la presi\'on justo en la base del tubo (el agua empuja el pist\'on
hacia afuera).
\end{dcl}
Ninguna otra fuerza act\'ua: no hay roce (el pist\'on ``puede desplazarse libremente''),
y como el pist\'on est\'a en reposo, estas tres fuerzas se equilibran exactamente.

\textbf{Equilibrio del pist\'on} (\'area $A=\pi D^2/4$): la presi\'on atmosf\'erica
$P_{atm}$ empuja el pist\'on hacia arriba (hacia el tubo); el peso $Mg$ y la presi\'on
del agua $P_b$ (justo sobre el pist\'on, en la \textbf{base} del tubo) empujan hacia
abajo:
\begin{equation}
Mg+P_b\,A=P_{atm}\,A
\quad\Longrightarrow\quad
P_b=P_{atm}-\frac{Mg}{A}
\end{equation}
\textbf{Presi\'on en el tope} (extremo cerrado, arriba): bajando hidrost\'aticamente
$h$ desde la base hasta el tope, la presi\'on \textbf{disminuye} (el tope est\'a m\'as
arriba):
\begin{equation}
P_t=P_b-\rho g h=P_{atm}-\frac{Mg}{A}-\rho g h
\end{equation}
La diferencia de presi\'on que debe resistir el tubo (entre el interior en el tope y el
exterior atmosf\'erico) es
\begin{equation}
|\Delta P_{tubo}|=\frac{Mg}{A}+\rho g h
\end{equation}
Expresada en metros de columna de agua (mca), dividiendo por $\rho g$:
\begin{equation}
\frac{1}{\rho g}\left(\frac{Mg}{A}+\rho g h\right)
\end{equation}
Con $M=65{,}6$ g, $h=19{,}5$ cm, $D=2$ cm ($A=\pi(0{,}01)^2=3{,}1416\times10^{-4}$
m$^2$):
\begin{equation}
\frac{Mg}{A}=2048{,}4\ \text{Pa}\ ,\qquad \rho g h=1912{,}95\ \text{Pa}
\end{equation}
\begin{center}
\fbox{\parbox{0.85\linewidth}{\centering
$\dfrac{1}{\rho g}\left(\dfrac{Mg}{A}+\rho gh\right)\approx0{,}404\ \text{mca}$\\[2pt]
Como $0{,}404\ \text{mca}>\Delta P=0{,}31\ \text{mca}$ (el m\'aximo que resiste el
tubo), \textbf{el tubo se rompe}.}}
\end{center}

%======================================================================
\newpage
\noindent{\LARGE\bfseries Preguntas 3 y 4 --- Freno de disco viscoso}

\noindent Un disco de di\'ametro $D$ gira entre dos paredes fijas, separado de cada una
por un lubricante de viscosidad $\mu$ y espesores $e_1$, $e_2$.

\begin{supuestos}
\item Flujo de Couette en cada espacio: el perfil de velocidad es lineal a trav\'es del
espesor (una pared fija a velocidad 0, el disco a velocidad $\omega r$).
\item Los dos lados del disco son independientes: cada uno aporta su propio torque
resistente, y se \textbf{suman}.
\end{supuestos}

\begin{dcl}
\item \textbf{Torque externo aplicado $T$} --- el que hay que ejercer para mantener
$\omega_0$ constante (s\'olo en P3; en P4 se deja de aplicar).
\item \textbf{Torque viscoso resistente del lado 1} --- el lubricante de espesor $e_1$
frena al disco, en sentido \textbf{opuesto} a la rotaci\'on.
\item \textbf{Torque viscoso resistente del lado 2} --- an\'alogo, con espesor $e_2$.
\end{dcl}
No hay fuerzas ``lineales'' que listar aqu\'i (el disco no se traslada, s\'olo gira): el
equilibrio es puramente rotacional, y son estos dos torques viscosos los que hacen todo
el trabajo --- en P3 equilibran al torque aplicado; en P4, al no haber torque aplicado,
son ellos solos los que desaceleran el disco.

\textbf{Parte P3 --- Torque para girar a $\omega_0$ constante.}
En un radio $r$, la velocidad del disco es $\omega r$; como el perfil es lineal a
trav\'es del espesor $e$ (un lado fijo en $0$, el otro en $\omega r$), el esfuerzo de
corte es
\begin{equation}
\tau(r)=\mu\,\frac{\omega r}{e}
\end{equation}
Este esfuerzo act\'ua sobre un anillo de \'area $dA=2\pi r\,dr$ a radio $r$, generando
una fuerza tangencial $dF=\tau\,dA$ y, por lo tanto, un torque elemental (brazo de
palanca $r$):
\begin{equation}
dT=r\,dF=r\cdot\mu\frac{\omega r}{e}\cdot2\pi r\,dr=\frac{2\pi\mu\omega}{e}\,r^3\,dr
\end{equation}
Integrando sobre todo el disco, de $r=0$ a $r=R$:
\begin{equation}
T_{\text{un lado}}=\frac{2\pi\mu\omega}{e}\int_0^R r^3\,dr
=\frac{2\pi\mu\omega}{e}\cdot\frac{R^4}{4}=\frac{\pi\mu\omega R^4}{2e}
\end{equation}
Como el disco tiene \textbf{dos} caras, cada una con su propio espesor, los torques se
suman:
\begin{equation}
\boxed{T=\frac{\pi\mu\omega R^4}{2}\left(\frac1{e_1}+\frac1{e_2}\right)}
\end{equation}
Con $e_1=0{,}25$ mm, $e_2=0{,}375$ mm, $\mu=302$ mPa\,s, $D=16$ cm ($R=0{,}08$ m),
$\omega_0=142$ rpm $=142\cdot\frac{2\pi}{60}=14{,}87$ rad/s:
\begin{center}
\fbox{$T\approx1{,}93\ \text{N}\cdot\text{m}$}
\end{center}

\textbf{Parte P4 --- Tiempo para llegar a $\omega_0/2$ (aqu\'i est\'a la integraci\'on).}
Al dejar de aplicar torque externo en $t=0$, el disco desacelera por el torque viscoso
resistente. La segunda ley de Newton para rotaci\'on ($I\,d\omega/dt=$ torque neto) da:
\begin{equation}
I\,\frac{d\omega}{dt}=-T=-\alpha\,\omega
\ ,\qquad
\alpha\equiv\frac{\pi\mu R^4}{2}\left(\frac1{e_1}+\frac1{e_2}\right)
\end{equation}
(usamos $T=\alpha\omega$ porque el torque viscoso es \textbf{proporcional} a $\omega$,
seg\'un la f\'ormula de P3 --- $\alpha$ es exactamente el mismo $T$ de P3 pero dividido
por $\omega$). Esta es una ecuaci\'on diferencial de \textbf{variables separables}:
\begin{equation}
I\,\frac{d\omega}{dt}=-\alpha\omega
\quad\Longrightarrow\quad
\frac{d\omega}{\omega}=-\frac{\alpha}{I}\,dt
\end{equation}
Integrando ambos lados (izquierda en $\omega$, desde $\omega_0$ hasta $\omega$; derecha
en $t$, desde $0$ hasta $t$):
\begin{equation}
\int_{\omega_0}^{\omega}\frac{d\omega'}{\omega'}=-\frac{\alpha}{I}\int_0^t dt'
\quad\Longrightarrow\quad
\ln\left(\frac{\omega}{\omega_0}\right)=-\frac{\alpha}{I}\,t
\end{equation}
\begin{equation}
\boxed{\omega(t)=\omega_0\,\exp\!\left(-\frac{\alpha}{I}\,t\right)}
\end{equation}
Es un decaimiento exponencial cl\'asico (misma estructura matem\'atica que un circuito
RC o un decaimiento radiactivo: la tasa de cambio es proporcional a la cantidad misma).
Usando el momento de inercia dado, $I=MR^2/2$, y pidiendo $\omega=\omega_0/2$:
\begin{equation}
\frac{\omega_0}{2}=\omega_0\exp\!\left(-\frac{\alpha}{I}t\right)
\quad\Longrightarrow\quad
\ln\!\left(\frac12\right)=-\frac{\alpha}{I}t
\quad\Longrightarrow\quad
\boxed{t=\frac{I\ln2}{\alpha}}
\end{equation}
Con $M=14$ g, $R=0{,}08$ m ($I=4{,}48\times10^{-5}\ \text{kg}\cdot\text{m}^2$) y
$\alpha=0{,}1295\ \text{N}\cdot\text{m}\cdot\text{s}$ (el mismo coeficiente de P3,
evaluado num\'ericamente sin el factor $\omega$):
\begin{center}
\fbox{$t\approx0{,}00024\ \text{s}=0{,}24\ \text{ms}$}
\end{center}

%======================================================================
\newpage
\noindent{\LARGE\bfseries Preguntas 5 y 6 --- Gas rotando a alta presi\'on}

\noindent Un recipiente en forma de sector circular (\'angulo $\theta$, radio $R$, alto
$e$) gira a velocidad angular $\omega$ constante, lleno de gas isot\'ermico (constante
espec\'ifica $R_0$, temperatura $T$). Se desprecia $\vec g$.

\begin{supuestos}
\item Equilibrio en el sistema rotante: el gradiente de presi\'on radial equilibra
exactamente la aceleraci\'on centr\'ipeta, $dP/dr=\rho\omega^2r$.
\item Gas ideal, isot\'ermico: $P=\rho R_0T$ con $T$ constante en todo el proceso.
\end{supuestos}

\textbf{Parte P5 --- Radio m\'inimo.} Sustituyendo $\rho=P/(R_0T)$ en el equilibrio
radial:
\begin{equation}
\frac{dP}{dr}=\frac{P}{R_0T}\,\omega^2r
\quad\Longrightarrow\quad
\frac{dP}{P}=\frac{\omega^2}{R_0T}\,r\,dr
\end{equation}
Integrando desde el centro ($r=0$, $P=P_m$, la presi\'on m\'inima que puede alcanzar el
gas) hasta el borde ($r=R$, $P=P_d$, la presi\'on de dise\~no):
\begin{equation}
\ln\!\left(\frac{P_d}{P_m}\right)=\frac{\omega^2R^2}{2R_0T}
\quad\Longrightarrow\quad
\boxed{R=\sqrt{\frac{2R_0T}{\omega^2}\ln\!\left(\frac{P_d}{P_m}\right)}}
\end{equation}
Con $\theta=23^\circ$, $P_m=0{,}0045$ atm, $P_d=17{,}4$ atm, $R_0=213$ J/(kg\,K),
$T=293$ K, $\omega=3248$ rpm $=340{,}13$ rad/s:
\begin{center}
\fbox{$R\approx2{,}99\ \text{m}\approx3\ \text{m}$}
\end{center}

\textbf{Parte P6 --- \textquestiondown Es viable el experimento?} La presi\'on en
cualquier radio, integrando desde el centro, es $P(r)=P_m\exp\!\left(\frac{\omega^2
r^2}{2R_0T}\right)$. La masa de gas atrapada en el sector, \textbf{despu\'es} de girar,
es
\begin{equation}
M=\int_0^\theta\!\!\int_0^R \rho(r)\,e\,r\,dr\,d\varphi
=\theta e\int_0^R\frac{P(r)}{R_0T}\,r\,dr
=\frac{\theta e\,P_m}{R_0T}\int_0^R\exp\!\left(\frac{\omega^2}{2R_0T}r^2\right)r\,dr
\end{equation}
Usando la ayuda del enunciado ($\int_0^R\exp(Ar^2)r\,dr=\frac1{2A}[\exp(AR^2)-1]$, con
$A=\omega^2/(2R_0T)$):
\begin{equation}
M=\frac{\theta e\,P_m}{R_0T}\cdot\frac{R_0T}{\omega^2}\left[\exp\!\left(\frac{\omega^2
R^2}{2R_0T}\right)-1\right]=\frac{\theta e\,P_m}{\omega^2}\left(\frac{P_d}{P_m}-1\right)
\end{equation}
(se us\'o $\exp(\omega^2R^2/2R_0T)=P_d/P_m$, el resultado de P5). Esta masa es la
\textbf{misma} que hab\'ia \textbf{antes} de girar, cuando el gas llenaba el mismo
sector con densidad uniforme $\rho_0$ (y presi\'on $P_0=\rho_0R_0T$):
\begin{equation}
M_{\text{antes}}=\rho_0\cdot\underbrace{\frac{\theta R^2}{2}}_{\text{\'area sector}}\cdot e
\end{equation}
Igualando ambas masas y despejando $\rho_0$ (equivalentemente $P_0$):
\begin{equation}
\boxed{\rho_0=\frac{2P_m}{\omega^2R^2}\left(\exp\!\left(\frac{\omega^2R^2}{2R_0T}
\right)-1\right)}
\end{equation}
Num\'ericamente: $\rho_m=P_m/(R_0T)=0{,}0073\ \text{kg/m}^3$, $\rho_0\approx3{,}42\
\text{kg/m}^3$, y por gas ideal $P_0=\rho_0R_0T\approx213\,386$ Pa $\approx2{,}11$ atm.
\begin{center}
\fbox{\parbox{0.85\linewidth}{\centering
$P_0\approx2{,}11\ \text{atm}>P_x=2\ \text{atm}$\\[2pt]
Se necesita m\'as presi\'on inicial que la m\'axima disponible experimentalmente:
\textbf{no es viable}.}}
\end{center}

%======================================================================
\newpage
\noindent{\LARGE\bfseries Preguntas 7 y 8 --- Compuerta semicil\'indrica entre dos
estanques}

\noindent Dos estanques comparten una pared con una compuerta semicil\'indrica s\'olida
(radio $R$, profundidad $L=1$ m, masa $M$), \textbf{articulada exactamente en su centro
de curvatura}. El estanque 1 tiene gas a presi\'on absoluta $P_g$ y agua hasta $h_1$; el
estanque 2 est\'a abierto a la atm\'osfera con agua hasta $h_2$.

\begin{supuestos}
\item La r\'otula est\'a en el centro de curvatura --- \'esta es la diferencia clave
respecto de la compuerta semicil\'indrica de la Clase Auxiliar (ah\'i estaba en el
extremo).
\item El estanque 1 moja directamente s\'olo el cuarto inferior de la cara plana de la
compuerta (altura $R$, desde el piso hasta la r\'otula).
\item La superficie curva queda expuesta al estanque 2, que la moja completa (altura
$2R$, con la r\'otula exactamente en su punto medio).
\item Densidad del agua $\rho=1000$ kg/m$^3$ en ambos estanques.
\end{supuestos}

\begin{dcl}
\item \textbf{Peso propio $Mg$} --- hacia abajo, aplicado en el centroide de la
compuerta ($4R/3\pi$ de la r\'otula, f\'ormula dada en el enunciado).
\item \textbf{Presi\'on del estanque 1} --- s\'olo sobre el cuarto inferior de la
compuerta (la \'unica parte en contacto con esa agua); resultante horizontal $F_2$,
aplicada bajo la r\'otula (a profundidad $y_R>x_G$, como en cualquier compuerta plana).
\item \textbf{Presi\'on del estanque 2} --- sobre la superficie curva completa; tiene
componente horizontal (fuerza sobre la proyecci\'on de altura $2R$) \textbf{y}
componente vertical (empuje $E$, en el centroide del semic\'irculo) --- pero, como se
muestra m\'as abajo, su torque \textbf{neto} respecto de la r\'otula es cero.
\item \textbf{Reacci\'on de la r\'otula} --- las dos componentes de fuerza que el
pasador ejerce sobre la compuerta; por ser una articulaci\'on, \textbf{no} aporta
torque respecto de s\'i misma (por eso no aparece en el balance de momentos, aunque s\'i
importar\'ia si pidieran esa reacci\'on).
\end{dcl}
Estas cuatro fuerzas son las que entran al balance de momentos. El enunciado menciona
adem\'as una posible ``reacci\'on por estar apoyada a la pared'' --- en la pauta del
profesor esto se resuelve sin necesitar una fuerza de apoyo aparte: la condici\'on
``ni rota ni hay reacci\'on'' es exactamente la masa $M$ que deja el balance de
momentos en cero usando s\'olo estas cuatro fuerzas (la r\'otula, al ser articulaci\'on,
no aporta torque; si el balance no cerrara, ah\'i s\'i har\'ia falta una reacci\'on
extra).

\textbf{Dos herramientas que se usan repetidamente en este problema:}

\textit{(1) Fuerza y torque horizontales sobre una superficie curva = los de su
proyecci\'on plana vertical.} Para cualquier superficie (plana o curva), la componente
\textbf{horizontal} de la fuerza de presi\'on, y su torque respecto de un eje
horizontal, son \textbf{exactamente iguales} a los de una compuerta plana vertical del
mismo alto y ancho --- porque $dF_x=P\,dA\cos\theta=P\,dA_{proy}$ es, punto a punto, la
misma integral que para una placa plana. Por eso se puede usar directamente
$y_R=y_G+I_{xG}/(Sy_G)$.

\textit{(2) Torque nulo de una superficie que rodea completamente el eje de giro.} Si
un fluido moja \textbf{toda} una superficie cil\'indrica centrada exactamente en la
r\'otula, la presi\'on en cada punto apunta en direcci\'on radial (pasa por el centro),
as\'i que su torque total (horizontal $+$ vertical) respecto de la r\'otula es
\textbf{cero} --- la misma idea de la Clase Auxiliar, pero aqu\'i s\'i se cumple.

\textbf{Fuerza y torque del estanque 1 (cuarto inferior, altura $R$).}
Usando la herramienta (1), con $A=R\cdot L$ y $x_G=h_1-R/2$ (profundidad del centroide
de esa franja bajo la superficie libre):
\begin{equation}
\boxed{F_2=P_g\,A+\rho g\left(h_1-\frac R2\right)A}
\end{equation}
Para el torque respecto de la r\'otula ($O=h_1-R$, su profundidad), se usa la
descomposici\'on cl\'asica de $y_R$: el t\'ermino de la presi\'on manom\'etrica
($\rho g$) aporta $I_{xG}+x_G(x_G-O)A$, y el t\'ermino de presi\'on uniforme del gas
$P_g$ aporta su propio brazo, $R/2$ (el punto medio de esa franja de altura $R$):
\begin{equation}
\boxed{T_2=\rho g\Big[I_{xG}+x_G(x_G-O)A\Big]+P_g\cdot\frac R2\cdot R}
\ ,\qquad I_{xG}=\frac{L R^3}{12}
\end{equation}
Con $h_1=1{,}26$ m, $R=0{,}56$ m, $P_g=1{,}5$ atm:
\begin{center}
\fbox{$F_2\approx94\,727\ \text{N}\ ,\qquad T_2\approx25\,357\ \text{N}\cdot\text{m}$}
\end{center}
Como el estanque 1 s\'olo moja un cuarto de la compuerta (\textbf{no} toda la
superficie cil\'indrica alrededor de la r\'otula), \textbf{no} aplica la herramienta
(2): este $T_2$ es un torque neto real, no nulo.

\textbf{Verificaci\'on de la herramienta (2), usando el estanque 2 (altura completa
$2R$).} El empuje (fuerza vertical) sobre un semicilindro completo, aplicado en el
centroide del semic\'irculo ($4R/3\pi$ del centro, dato del enunciado):
\begin{equation}
E=\rho g\,\frac{\pi R^2}{2}\,L
\ ,\qquad
T_E=E\cdot\frac{4R}{3\pi}=\rho g L\,\frac{2R^3}{3}
\end{equation}
Para la fuerza horizontal, con altura completa $2R$ y la r\'otula exactamente en el
centro de esa franja ($x_G=O=h_2-R$), el t\'ermino $x_G(x_G-O)$ se anula, dejando s\'olo
el t\'ermino de forma:
\begin{equation}
T_2^{(2R)}=\rho g\,I_{xG}^{(2R)}=\rho g\,\frac{L(2R)^3}{12}=\rho g L\,\frac{2R^3}{3}
\end{equation}
\textbf{Son id\'enticos:} $T_E=T_2^{(2R)}$, para cualquier $R$, $\rho$, $L$ --- \'esta
es exactamente la herramienta (2) confirmada algebraicamente: el torque vertical y el
horizontal de una superficie que rodea completamente la r\'otula se cancelan
exactamente, dando torque neto cero del estanque 2.

\textbf{Parte P8 --- masa l\'imite.} Como el estanque 2 no aporta torque neto (recién
verificado) y la superficie plana embebida (herramienta del peso) no aporta empuje
(es plana), la \'unica forma de que la compuerta est\'e en equilibrio \textbf{sin}
reacci\'on del apoyo es que el peso equilibre exactamente el \'unico torque neto que
queda: el del estanque 1.
\begin{equation}
\sum M_{r\acute otula}=0:\qquad Mg\cdot\frac{4R}{3\pi}=T_2
\end{equation}
\begin{equation}
\boxed{M=\frac{3\pi T_2}{4gR}}
\end{equation}
Con $T_2\approx25\,357$ N$\cdot$m, $R=0{,}56$ m:
\begin{center}
\fbox{$M\approx10\,875\ \text{kg}$}
\end{center}

\vspace{4pt}
\noindent{\small\textit{Nota:} la parte m\'as sutil de este problema es notar que el
"torque cero" (herramienta 2) s\'olo se aplica al estanque 2, porque \'este s\'i moja
la compuerta por completo alrededor de la r\'otula. El estanque 1, al mojar s\'olo un
cuarto, \textbf{s\'i} genera torque neto --- y es exactamente ese torque el que el peso
debe equilibrar.}

\end{document}
